\documentclass[a4paper,12pt]{article}
\usepackage[T2A]{fontenc}
\usepackage[utf8]{inputenc}
\usepackage[russian]{babel}
\usepackage{amsmath,amssymb,mathtools}
\usepackage{helvet}
\renewcommand{\familydefault}{\sfdefault}
\usepackage{icomma}
\usepackage{geometry}
\geometry{left=21mm,right=21mm,top=20mm,bottom=21mm}
\usepackage{microtype}
\usepackage{tikz}
\usetikzlibrary{arrows.meta,positioning,shapes.geometric}
\usepackage{pgfplots}
\pgfplotsset{compat=1.18}
\usepackage{booktabs,tabularx,longtable}
\usepackage{enumitem}
\usepackage{fancyhdr}
\usepackage{parskip}
\usepackage{setspace}
\usepackage{xcolor}
\usepackage{hyperref}
\usepackage{needspace}
\definecolor{accent}{HTML}{008080}
\definecolor{darkaccent}{HTML}{004D4D}
\definecolor{textgray}{HTML}{333333}
\definecolor{lightgray}{HTML}{707A7A}
\hypersetup{colorlinks=true,linkcolor=darkaccent,urlcolor=darkaccent}
\color{textgray}
\setstretch{1.14}
\setlength{\parindent}{0pt}
\setlength{\parskip}{6pt plus 2pt minus 1pt}
\setlist[itemize]{leftmargin=1.3em,itemsep=2pt,topsep=3pt}
\setlist[enumerate]{leftmargin=2.0em,itemsep=12pt,topsep=5pt}
\pagestyle{fancy}
\fancyhf{}
\fancyhead[L]{\footnotesize\color{lightgray}Неравенства и числовые промежутки}
\fancyhead[R]{\footnotesize\color{lightgray}Екатерина Скелли}
\fancyfoot[C]{\footnotesize\color{lightgray}\thepage}
\renewcommand{\headrulewidth}{0pt}
\renewcommand{\footrulewidth}{0pt}
\setlength{\headheight}{15pt}
\newcommand{\chap}[1]{\needspace{7\baselineskip}\section*{\color{darkaccent}#1}}
\newcommand{\ex}[1]{\par\smallskip\textcolor{darkaccent}{\textbf{Пример.}} #1\par}
\newcommand{\key}[1]{\textcolor{darkaccent}{\textbf{#1}}}
\begin{document}
\thispagestyle{empty}
\begin{center}
\vspace*{27mm}
{\large\color{accent}\textbf{МАТЕМАТИКА}}\par
\vspace{28mm}
{\fontsize{28}{34}\selectfont\color{darkaccent}\textbf{Чек-лист}}\par
\vspace{9mm}
{\Large Неравенства, числовые промежутки\\[2mm]и графический метод}\par
\vspace{12mm}
{\normalsize Подготовка к контрольной работе}\par
\vspace{25mm}
{\large\textbf{Екатерина Скелли}}\par
\vspace{5mm}
4 ноября 2026 г.
\end{center}
\vfill
\begin{center}
{\small Лёвин Артём Александрович}\\[2pt]
{\small эксклюзивно для Екатерины Скелли}
\end{center}
\vspace{10mm}
\begin{center}
\begin{tikzpicture}
\draw[accent!65,line width=0.75pt] (-4.4,0) .. controls (-2.1,0.35) and (-1.1,-0.35) .. (0,0)
 .. controls (1.15,0.34) and (2.15,-0.3) .. (4.4,0);
\end{tikzpicture}
\end{center}
\clearpage

\chap{1. Что необходимо уметь}
\begin{itemize}[label=$\square$]
\item Сравнивать обыкновенные и десятичные дроби, включая отрицательные числа.
\item Переходить между неравенством, числовым промежутком и изображением на оси.
\item Находить пересечение и объединение числовых множеств.
\item Решать линейные и двойные неравенства, учитывать знак при делении.
\item Оценивать суммы, разности, произведения, частные и степени переменных.
\item Читать линейные графики и находить условия их пересечения.
\item Разбирать случаи, когда коэффициент зависит от параметра.
\end{itemize}

\chap{2. Линейные неравенства}
К одному и тому же выражению можно прибавить любое число или вычесть его из обеих частей неравенства. При умножении или делении на \emph{положительное} число знак сохраняется; на \emph{отрицательное} число знак меняется.
\[
-3x>12\quad\Longleftrightarrow\quad x<-4.
\]
Перед делением на выражение с параметром сначала определите его знак. \key{На ноль делить нельзя.}

\ex{Решим $3(2x-1)-2(x+4)\ge x-5$:
\[
6x-3-2x-8\ge x-5\ \Longleftrightarrow\ 3x\ge6
\ \Longleftrightarrow\ x\ge2.
\]}

\chap{3. Промежутки и множества}
\[
\begin{array}{c@{\qquad}c}
x>2\ \Longleftrightarrow\ x\in(2;+\infty), &
x\ge2\ \Longleftrightarrow\ x\in[2;+\infty),\\[4pt]
-3<x\le4\ \Longleftrightarrow\ x\in(-3;4]. &
\end{array}
\]
Круглая скобка означает, что граница исключена; квадратная --- включена. При $\pm\infty$ всегда ставится круглая скобка. На числовой прямой исключённая граница --- пустая точка, включённая --- закрашенная.

\key{Пересечение} $A\cap B$ --- общая часть множеств. \key{Объединение} $A\cup B$ --- все точки, принадлежащие хотя бы одному из них.

\ex{Если $A=(-4;2]$, $B=[0;5)$, то
\[
A\cap B=[0;2],\qquad A\cup B=(-4;5).
\]}
\clearpage

\chap{4. Двойные неравенства и оценки выражений}
Действие нужно совершать \emph{со всеми тремя частями} двойного неравенства. Если умножаем на отрицательное число, обе стрелки меняют направление; границы затем записывают по возрастанию.

\ex{Если $-3<a<2$, оценим $5-2a$:
\[
-3<a<2\ \Longrightarrow\ -4<-2a<6
\ \Longrightarrow\ \boxed{1<5-2a<11}.
\]}

Для независимых переменных $a\in[a_1;a_2]$, $b\in[b_1;b_2]$:
\begin{itemize}
\item сумму оцениваем сложением соответствующих границ;
\item разность: наименьшее значение $a_1-b_2$, наибольшее $a_2-b_1$;
\item произведение: сравниваем четыре числа $a_1b_1$, $a_1b_2$, $a_2b_1$, $a_2b_2$;
\item частное: сначала убеждаемся, что $0\notin[b_1;b_2]$, затем сравниваем четыре частных граничных значений.
\end{itemize}
Для открытых промежутков найденные граничные значения могут быть лишь пределами: скобки будут круглыми. Для квадратов и обратных величин учитывайте также знак, ноль и возможные внутренние значения.

\ex{При $-5<x<-2$ верно
\[
4<x^2<25,\qquad -\frac12<\frac1x<-\frac15.
\]}

\textcolor{accent}{\textbf{Важное ограничение.}} Если знаменатель может равняться нулю, выражение в этой точке не определено, а оценка частного через четыре граничных значения может быть неверной. Например, при $-3<a<1$ и $-2<b<-1$ знаменатель $2a-3b$ способен обращаться в ноль. Поэтому для дроби $\dfrac{a+b}{2a-3b}$ нельзя автоматически записывать конечные границы по углам прямоугольника допустимых значений.

\chap{5. Графический смысл неравенства}
Если задано $y=f(x)$, условие $f(x)>0$ требует найти те $x$, при которых график расположен \emph{выше} оси $Ox$. Условие $f(x)\ge c$ означает, что график находится на уровне $y=c$ или выше. Для строгого знака точки равенства исключают.

\ex{Для $y=2x+4$ график --- прямая через точки $(0;4)$ и $(-2;0)$. Условие $2x+4\ge0$ выполняется при $x\ge-2$.}
\begin{center}
\begin{tikzpicture}
\begin{axis}[width=.72\linewidth,height=6.2cm,axis lines=middle,xmin=-4,xmax=2.5,ymin=-4,ymax=8,xtick={-4,-3,-2,-1,0,1,2},ytick={-4,-2,0,2,4,6,8},grid=major,grid style={line width=.2pt,draw=black!9},axis line style={->},tick label style={font=\scriptsize},xlabel={$x$},ylabel={$y$},xlabel style={at={(axis description cs:1,.32)},anchor=west},ylabel style={at={(axis description cs:.65,1)},anchor=south},clip=false]
\addplot[domain=-4:2.0,samples=2,color=accent,very thick]{2*x+4};
\addplot[only marks,mark=*,mark size=2.2pt,color=darkaccent]coordinates{(-2,0)(0,4)};
\node[font=\footnotesize,above left] at (axis cs:-2,0){$(-2;0)$};
\node[font=\footnotesize,right] at (axis cs:0,4){$(0;4)$};
\end{axis}
\end{tikzpicture}
\end{center}
\clearpage

\chap{6. Пересечение двух прямых}
Если $y=f(x)$ и $y=g(x)$ пересекаются, то в общей точке их ординаты равны: $f(x)=g(x)$. Сначала находим абсциссу $x$, затем подставляем её в любое из двух уравнений для вычисления $y$.
\[
2x-3=-5x+2\quad\Longrightarrow\quad 7x=5\quad\Longrightarrow\quad x=\frac57,
\]
\[
y=2\cdot\frac57-3=-\frac{11}{7}.
\]
Абсцисса --- координата $x$; ордината --- координата $y$. Если в формуле присутствует параметр $a$, сначала можно выразить $x$ через $a$, затем применить условие, например $x<0$.

\chap{7. Неравенство с параметром}
Параметр --- число, от которого может зависеть вид множества решений. В неравенстве $kx\le b$ определяющий вопрос: \emph{каков знак коэффициента $k$?}
\[
\begin{array}{ll}
k>0:&x\le\dfrac bk,\\[7pt]
k<0:&x\ge\dfrac bk,\\[7pt]
k=0:&0\le b\text{ --- проверяем числовое утверждение.}
\end{array}
\]
Если при $k=0$ утверждение верно, подходят все действительные $x$; если неверно, решений нет.

\ex{В неравенстве $(3-a)x\le2a-1$ при $a=3$ получаем $0\le5$, поэтому подходит любое $x$. При $a<3$ знак сохраняется, а при $a>3$ меняется.}

\chap{8. Доказательство неравенств}
Если $0<a<b$, то $b+1>0$ и $b+2>0$. Следовательно, при умножении на их произведение знак сохраняется:
\[
\frac{a+1}{b+1}<\frac{a+2}{b+2}
\quad\Longleftrightarrow\quad
(a+1)(b+2)<(a+2)(b+1)
\]
\[
\Longleftrightarrow\quad ab+2a+b+2<ab+a+2b+2
\quad\Longleftrightarrow\quad a<b.
\]
Последнее верно по условию, что и доказывает исходное неравенство.

\clearpage

\begin{center}
{\Large\color{darkaccent}\textbf{Алгоритм: линейное неравенство с параметром}}\par
\vspace{8mm}
\end{center}
\vspace{16mm}
\begin{center}
\begin{tikzpicture}[>=Latex, every node/.style={font=\small},
 start/.style={ellipse,draw=accent,minimum width=38mm,minimum height=11mm,align=center},
 action/.style={rectangle,rounded corners=2pt,draw=accent,minimum width=44mm,minimum height=14mm,align=center,text width=42mm},
 decision/.style={diamond,draw=accent,aspect=2.0,align=center,inner sep=3pt,text width=27mm},
 arr/.style={->,draw=darkaccent,line width=.9pt}]
\node[start] (start) at (0,0) {$kx\le b$};
\node[decision] (pos) at (0,-2.3) {$k>0$?};
\node[action] (plus) at (5,-2.3) {$x\le\dfrac bk$\\знак сохраняется};
\node[decision] (neg) at (0,-5.5) {$k<0$?};
\node[action] (minus) at (5,-5.5) {$x\ge\dfrac bk$\\знак меняется};
\node[action] (zero) at (0,-8.4) {$k=0$: проверить $0\le b$};
\node[action] (yes) at (-3,-11.2) {Верно: $x\in\mathbb R$};
\node[action] (no) at (3,-11.2) {Неверно: $\varnothing$};
\draw[arr] (start)--(pos);
\draw[arr] (pos)--node[above]{да}(plus);
\draw[arr] (pos)--node[left]{нет}(neg);
\draw[arr] (neg)--node[above]{да}(minus);
\draw[arr] (neg)--node[left]{нет}(zero);
\draw[arr] (zero.south)--++(0,-.55)-|node[near end,left]{да}(yes.north);
\draw[arr] (zero.south)--++(0,-.55)-|node[near end,right]{нет}(no.north);
\end{tikzpicture}
\end{center}
\clearpage

\section*{\color{darkaccent}Тренировочный блок: 20 упражнений}
Выполняйте задания по порядку. Для неравенств записывайте множество решений, для оценок --- точные границы. В заданиях с графиками приводите чертёж. Ответы находятся после тренировочного блока.

\subsection*{\color{accent}Базовый уровень}
\begin{enumerate}[start=1]
\item Сравните числа, поставив знак $<$, $>$ или $=$:
\[
\text{а) }0{,}64\ \square\ \frac9{14};\qquad
\text{б) }-\frac{11}{8}\ \square\ -1{,}36.
\]
\item Запишите промежуток для неравенства $-4<x\le3$ и изобразите его на числовой прямой.
\item Запишите через неравенства: \textbf{а)} $x\in[2;+\infty)$; \textbf{б)} $x\in(-\infty;-1)$.
\item Даны $A=(-5;2]$, $B=[-1;7)$. Найдите $A\cap B$ и $A\cup B$.
\item Даны $A=(-\infty;1)$, $B=[1;6)$. Найдите $A\cap B$ и $A\cup B$.
\item Решите неравенство $4x-7<2x+9$.
\item Решите неравенство $3(2x-1)-2(x+4)\ge x-5$.
\item Решите двойное неравенство $-3<2x+1\le9$.
\end{enumerate}

\subsection*{\color{accent}Средний уровень}
\begin{enumerate}[start=9]
\item Решите неравенство $\dfrac{5-2x}{3}>x-1$ и перечислите все его \textbf{целые} решения из отрезка $[-1;4]$.
\item Для каждого значения $a$ определите множество решений неравенства $2y+3a<2y+7$: \textbf{а)} $a=-4$; \textbf{б)} $a=3$.
\item Тариф А стоит $150+32n$ рублей, тариф Б --- $70+48n$ рублей, где $n$ --- целое неотрицательное число часов. При каком \textbf{наименьшем} $n$ тариф А дешевле тарифа Б?
\item Если $-2<a<3$, оцените $4-3a$.
\item Если $-2\le a\le1$ и $3\le b\le6$, оцените \textbf{а)} $a+b$; \textbf{б)} $a-b$.
\item Если $-4\le a\le-2$ и $1\le b\le3$, оцените произведение $ab$.
\item Если $-6<t<-3$, оцените \textbf{а)} $t^2$; \textbf{б)} $\dfrac1t$.
\item Если $-1<x<2$ и $1<y<3$, оцените дробь $\dfrac{2x+4}{5-y}$. Покажите отдельно оценки числителя и знаменателя.
\end{enumerate}

\subsection*{\color{accent}Повышенный уровень}
\begin{enumerate}[start=17]
\item Постройте график $y=2x-6$ по двум точкам. По графику найдите $x$, при которых \textbf{а)} $y\ge0$; \textbf{б)} $y<4$. Отметьте пересечения с осями координат.
\item Найдите координаты точки пересечения прямых $y=3x-2$ и $y=-x+6$.
\item \textbf{а)} Прямые $y=3x-a+4$ и $y=5x-3a-6$ пересекаются в точке с отрицательной абсциссой. При каких $a$ это верно?\\
\textbf{б)} Решите при всех действительных $a$ неравенство $(3-a)x\le2a-1$.
\item Докажите, что при $0<a<b$ справедливо
\[
\frac{a+1}{b+1}<\frac{a+2}{b+2}.
\]
Укажите, почему при преобразованиях знак неравенства можно сохранить.
\end{enumerate}
\clearpage

\section*{\color{darkaccent}Ответы к тренировочному блоку}
\renewcommand{\arraystretch}{1.4}
\small
\begin{longtable}{@{}p{10mm}p{\dimexpr\linewidth-16mm\relax}@{}}
\toprule
\textbf{№} & \textbf{Ответ}\\
\midrule
\endfirsthead
\toprule \textbf{№} & \textbf{Ответ}\\ \midrule
\endhead
1 & а) $0{,}64<\dfrac9{14}$; б) $-\dfrac{11}{8}<-1{,}36$\\
2 & $(-4;3]$\\
3 & а) $x\ge2$; б) $x<-1$\\
4 & $A\cap B=[-1;2]$; $A\cup B=(-5;7)$\\
5 & $A\cap B=\varnothing$; $A\cup B=(-\infty;6)$\\
6 & $x<8$\\
7 & $x\ge2$\\
8 & $-2<x\le4$\\
9 & $x<\dfrac85$; целые решения: $-1,\,0,\,1$\\
10 & а) $y\in\mathbb R$; б) $\varnothing$\\
11 & $n=6$\\
12 & $-5<4-3a<10$\\
13 & а) $1\le a+b\le7$; б) $-8\le a-b\le-2$\\
14 & $-12\le ab\le-2$\\
15 & а) $9<t^2<36$; б) $-\dfrac13<\dfrac1t<-\dfrac16$\\
16 & $2<2x+4<8$; $2<5-y<4$; $\dfrac12<\dfrac{2x+4}{5-y}<4$\\
17 & Точки пересечения с осями: $(3;0)$, $(0;-6)$; а) $x\ge3$; б) $x<5$\\
18 & $(2;4)$\\
19 & а) $a<-5$; б) при $a<3$: $x\le\dfrac{2a-1}{3-a}$; при $a=3$: $x\in\mathbb R$; при $a>3$: $x\ge\dfrac{2a-1}{3-a}$\\
20 & Умножаем на положительное $(b+1)(b+2)$; получаем равносильное $a<b$, верное по условию.\\
\bottomrule
\end{longtable}
\normalsize
\vfill
\begin{center}
\small\color{lightgray}Лёвин Артём Александрович --- эксклюзивно для Екатерины Скелли
\end{center}
\end{document}
